\section{Endpoint Reduction for the Height Parameter}
\label{app:endpoints}

The scalar comparison initially depends on two variables: the channel's
half-log-odds \(\logodds\) and one coordinate's comparison height \(v\).
This appendix removes the height variable analytically.
The reason is elementary: a suitable upper bound can decrease and then
increase, but cannot have an interior maximum. We only need its endpoints.

\subsection{Two Shape Properties}

Recall \(c(v)=\arcsin(\tanh v)^2/\tanh v\) and
\(F(v)=h(\tanh v)/\tanh v\). On \(v>0\),
\[
 c\text{ is concave},\qquad 1/F\text{ is convex}.
\]
Here is a short proof. Put
\(s(v)=\arcsin(\tanh v)/\sinh v\). Differentiating gives
\[
 c'(v)=2s(v)-s(v)^2,\qquad 0<s(v)<1,\qquad s'(v)<0.
\]
Indeed, \(\sinh v-\arcsin(\tanh v)\) starts at zero and has
positive derivative. Also
\[
 s'(v)=\frac{\tanh v-\arcsin(\tanh v)\cosh v}{\sinh^2v}<0,
\]
since \(\arcsin(\tanh v)>\tanh v>0\). Hence
\(c''(v)=2(1-s(v))s'(v)<0\), and \(0<c'(v)<1\).

For the reciprocal profile write
\[
 J(v)=\frac1{F(v)}=\frac{\sinh v}{{\mathcal H}(v)},\qquad
 {\mathcal H}(v)=\cosh v\,h(\tanh v)>0.
\]
The difference \(v\coth v-\log(2\cosh v)\) decreases to zero:
its derivative is \((\tanh v-v)/\sinh^2v<0\).
Consequently
\({\mathcal H}'(v)=\sinh v\log(2\cosh v)-v\cosh v<0\).
Direct differentiation now gives
\begin{align*}
 J'(v)&=\frac{\log(2\cosh v)}{{\mathcal H}(v)^2},\\
 J''(v)&=\frac{\tanh v\,{\mathcal H}(v)
 -2\log(2\cosh v){\mathcal H}'(v)}{{\mathcal H}(v)^3}>0.
\end{align*}

\subsection{Lower and Upper Endpoints}

Fix a channel with \(\ell=\logodds\ge7/2\), write \(r=\rho\),
and use the canonical multiplier \(p=(\lambda+1/\alpha)^2\).
The Parseval bound on the coordinate yield is
\[
 Y_P(v)=\frac{r c(v)+p[d_2/2+(d_1-d_2/2)F(\ell)/F(v)]}{v}.
\]
We prove below that \(Y_P\) has no interior maximum when
\(p d_2/(2r)\ge5/32\). Appendix~\ref{app:canonical-prerequisites}
verifies this condition on the whole stated half-line.
The LSI comparison covers heights through \(z=\ell/2\).
Define the last retained height \(V\) by
\(F(V)=F(\ell)/a_{\rm cut}\), where the local lemma covers
influences at least \(a_{\rm cut}\).
If \(V\le z\), the LSI bound covers everything. Otherwise the
only remaining checks are \(Y_P(z)<B/\ell\) and
\(Y_P(V)<B/\ell\).

The log-sum argument handles
the smaller channels without this coordinate comparison.

\subsection{A Uniform Rate for the Entropy Profile}

The cutoff proof needs no inverse formula for \(F\). It only needs
to know that \(F\) decreases fast enough. Define its relative decay
rate \(Q_F(v)=-F'(v)/F(v)\). We prove
\begin{equation}
                         Q_F(v)>8/5\qquad(v>0).
 \label{eq:analytic-profile-speed}
\end{equation}
Writing \(q_v=e^{-2v}\), differentiation of the entropy formula gives
\[
 Q_F(v)=\frac{4q_v[v+\log(1+q_v)]}
 {(1-q_v)[(1+q_v)\log(1+q_v)+2v q_v]}.
\]
After multiplication by positive denominators, the claimed bound is
equivalent to
\[
 v q_v(1+4q_v)+(5q_v-2+2q_v^2)\log(1+q_v)>0.
\]
If the logarithm has nonnegative coefficient, this is immediate.
Otherwise \(\log(1+q_v)<q_v\) reduces it to the stronger assertion
\(v-2+(4v+5)e^{-2v}>0\). This is immediate for \(v\ge2\).
For \(0\le v<2\), the ratio
\((2-v)e^{2v}/(4v+5)\) has derivative with the sign of
\(-8v^2+6v+7\). Its maximum is therefore at
\((3+\sqrt{65})/8\), between \(11/8\) and \(7/5\).
The maximum is less than \(5e^{14/5}/84<1\). The last fixed
exponential bound follows from the finite series and geometric
remainder in Section~\ref{app:fixed-arithmetic}. This proves the rate.

For the cutoff \(a_{\rm cut}=1/(1+(8/3)e^{-\ell})\), integration
between its actual height \(V\) and \(\ell\) now gives
\begin{align*}
 \frac85(\ell-V)&<\int_V^\ell Q_F(v)\,dv=-\log a_{\rm cut}\\
 &=\log\bigl(1+(8/3)e^{-\ell}\bigr)<\frac83e^{-\ell}.
\end{align*}
Thus \(0<\ell-V<(5/3)e^{-\ell}\), without evaluating an inverse
of \(F\).

\subsection{Endpoint Principle for the Original Yield}
\label{app:original-yield}

The original Parseval yield has no interior maximum once its constant
term pays for the following global curvature bound:
\begin{equation}
 \begin{split}
 &\frac{-c''(v)[vJ'(v)-J(v)]}{J''(v)}-[c(v)-vc'(v)]<\frac5{32},\\
 &\hspace{8em}J=1/F,\qquad v>0.
 \end{split}
 \label{eq:original-shape}
\end{equation}
We prove it below. Its use is simple. For positive constants
\(A_0,A_1\), the derivative of
\[
                    Y_P(v)=\frac{r c(v)+A_0+A_1 J(v)}v
\]
has the sign of
\[
 A_1-\frac{A_0+r[c(v)-vc'(v)]}{vJ'(v)-J(v)}.
\]
The denominator is positive. Differentiation and
\eqref{eq:original-shape} show that the fraction strictly decreases
if \(A_0/r\ge5/32\). The derivative of \(Y_P\) can therefore change
sign only from negative to positive: \(Y_P\) has no interior maximum.
This applies directly to the Parseval yield with
\(A_0=p d_2/2\) and \(A_1=p(d_1-d_2/2)F(\ell)\).
Appendix~\ref{app:canonical-prerequisites} proves
\(p d_2/(2r)>3/16>5/32\) for every \(\ell\ge7/2\).

\paragraph{Small-height entropy comparison.}
Recall \(L=\log2\). Continue to write
\({\mathcal H}(v)=\cosh v\,h(\tanh v)\), and put
\(\Lambda(v)=\log(2\cosh v)\). Recall that \({\mathcal H}'<0\) and
\(J'=\Lambda/{\mathcal H}^2\).
Put \(r_v=\tanh v\), \(\delta_v=1-r_v^2\). For a decreasing
function \(f\) and an increasing function \(g\) on \([0,1]\),
\begin{align*}
 &\int_0^1fg-\left(\int_0^1f\right)\left(\int_0^1g\right)\\
 &\quad=\frac12\int_0^1\!\int_0^1
 [f(s)-f(t)][g(s)-g(t)]\,ds\,dt\le0.
\end{align*}
Apply this to \(f(s)=1/(1+s)\), \(g(s)=1/(1-r_v^2s^2)\):
\[
 \frac{h(r_v)}{\delta_v}
 =\int_0^1\frac{ds}{(1+s)(1-r_v^2s^2)}
 \le L\int_0^1\frac{ds}{1-r_v^2s^2}=\frac{Lv}{r_v}.
\]
The integral identity follows by integrating the partial fractions
\begin{align*}
 \frac1{(1+s)(1-r^2s^2)}={}&
 \frac1{(1-r^2)(1+s)}\\
 &+\frac{r}{2(1+r)(1-rs)}\\
 &-\frac{r}{2(1-r)(1+rs)}.
\end{align*}

For \(0<v\le1\), this also gives \({\mathcal H}''<0\). Indeed,
\(v/\tanh v\) increases, so \(h(r_v)/\delta_v\le L\coth1<1\):
use \(e^2>7\), hence \(\coth1<4/3\), and \(L<3/4\). Differentiation gives
\({\mathcal H}''=(h(r_v)-\delta_v)/\sqrt{\delta_v}\).
Thus \({\mathcal H}^{-2}\) is increasing and convex on this interval.
The product rule gives \(J'''=(\Lambda {\mathcal H}^{-2})''>0\), since
\(\Lambda\) is positive, increasing, and convex. Hence
\begin{equation}
 vJ'(v)-J(v)=\int_0^v tJ''(t)\,dt
                  <\frac{v^2}{2}J''(v)\qquad(0<v\le1).
 \label{eq:original-profile-curvature}
\end{equation}

\paragraph{Accumulated curvature.}
We first strengthen the preceding integral bound to
\begin{equation}
 \frac{vJ'(v)-J(v)}{J''(v)}<\frac v2\qquad(v>0).
 \label{eq:profile-ratio-linear}
\end{equation}
For \(v\le1\), \eqref{eq:original-profile-curvature} already proves it.
For \(v\ge1\), put
\begin{align*}
 q&=e^{-2v},& T&=\frac{\log(1+q)}q,\\
 K&=2v+(1+q)T,& D&=vJ'-J.
\end{align*}
The entropy formula gives \(J=(1-q)/(qK)\), and differentiation gives
\begin{align*}
 qD&=1-\frac{1-q+2T}K+\frac{(1-q^2)T^2}{K^2},\\
 K'&=2T,\qquad 0<T'<q.
\end{align*}
Here primes denote derivatives in \(v\). The last bound follows from
the reciprocal trapezoid estimate: \(T'=2[T-1/(1+q)]<q\).
Also \(T\ge1-q/2\). The last numerator \((1-q^2)T^2\) increases,
so differentiating and dropping its positive derivative gives
\begin{align*}
 (qD)'\ge\frac1{K^2}\biggl[&2T(1-q+2T)-K(2q+2T')\\
 &-\frac{4(1-q^2)T^3}K\biggr]>0.
\end{align*}
The three terms are respectively greater than \(4\), less than
\(2\), and less than \(2\). Indeed, \(e^2>7\) gives
\(q<1/7\) and \(1-q,T>6/7\). Also
\(qK<e^{-2v}(2v+3/2)<1/2\), since this upper bound decreases
from a value below \((1/7)(7/2)\) at one. Finally \(K>2\) and \(T<1\).
Since \((qD)'=q(vJ''-2D)\), this proves
\eqref{eq:profile-ratio-linear}.

\paragraph{Scaling of angular curvature.}
Write
\[
 w(v)=\frac{\arcsin(\tanh v)}{\sinh v},
 \qquad k(v)=-c''(v)>0.
\]
The function \(w\) is decreasing and convex as a function of \(v^2\).
For a direct proof, put \(r_v=\tanh v\) and
\(\tau_v=\operatorname{sech}v\).
The elementary angle bound
\[
 \arcsin r_v\ge\frac{3r_v}{2+\tau_v},\qquad
 w(v)\ge\frac{3\tau_v}{2+\tau_v}
\]
follows by differentiation: the derivative of the first difference,
with respect to \(r_v\), is
\((1-\tau_v)^2/[\tau_v(2+\tau_v)^2]\ge0\).
Since \(3\tau_v/(2+\tau_v)>\tau_v^2\), the identity
\(w'=(\tau_v^2-w)/r_v\) also proves that \(w\) decreases.
Differentiate once more and substitute
this lower bound:
\[
 vw''-w'\ge
 \frac{\tau_v(1-\tau_v)}{r_v^2(2+\tau_v)}
 \bigl[(3+\tau_v)r_v+v\{4-(1+\tau_v)^3\}\bigr]>0.
\]
For the last sign, \(r_v\ge v\tau_v\) leaves
\(v(3-2\tau_v^2-\tau_v^3)>0\) inside the bracket.
The sign \(vw''-w'>0\) is exactly convexity in \(v^2\).

Consequently \(f(v)=1-w(v)\) is increasing and concave in \(v^2\),
and vanishes at zero. For \(0<t\le v\),
concavity gives \(f(t)\ge(t/v)^2f(v)\) and
\(f'(t)\ge(t/v)f'(v)\). Since \(k=2ff'\), integration yields
\begin{equation}
 c(v)-vc'(v)=\int_0^v t k(t)\,dt
                     \ge\frac{v^2}{5}k(v).
 \label{eq:angle-curvature-scaling}
\end{equation}
Also \(k(v)\le\tanh^3(v)/2<1/2\): in
\(k=2(1-w)(w-\delta)/r\), where \(r=\tanh v\) and
\(\delta=1-r^2\), maximize the quadratic in \(w\).
Combining the two estimates gives
\begin{align*}
 &\frac{k(v)[vJ'(v)-J(v)]}{J''(v)}-[c(v)-vc'(v)]\\
 &\qquad\le k(v)\left(\frac v2-\frac{v^2}5\right)<\frac5{32}.
\end{align*}
For \(v\le5/2\), use \(k<1/2\) and complete the square;
for larger \(v\), the bracket is nonpositive. This proves
\eqref{eq:original-shape}, and hence the endpoint principle.

